\documentclass{situationist-preprint}

\title{Possible Cases}
\msid{PID-10-b}
\msfamily{PID-10}
\msdiscipline{Personal Identity and the Self}
\mstype{Consequence}
\mskeywords{disembodied persons; closed emulation; locked-in syndrome; uploading; dualism; phase sortal}

\begin{document}
\maketitle

\begin{abstract}
PID-10-a argues that disembodied persons are possible in principle, by way of the closed emulation, and that no actual case is established. We derive which kinds of case qualify. A closed emulation, and an upload decoupled from all transducers and effectors, would be disembodied persons; a human being who is paralyzed and deprived of all sensation, a dreamer and a brain in a vat would not; a soul would be one only if substance dualism is true, which the result does not require. In no case does a human person become disembodied, though she may have a disembodied successor.
\end{abstract}

\section{The criterion}

On PID-10-a, a person is a thing with the capacities for rational thought and for thinking of itself as itself (Definition~2.1); a body of a thing is something through which it perceives or on which it acts directly (Assumption~2.6); and a thinking organism is its own body (Assumption~2.7). A thing is disembodied if nothing is its body (Definition~2.2). A case is therefore a case of a disembodied person exactly when something has the capacities of a person and perceives through nothing and acts directly on nothing.

\section{The cases}

\begin{proposition}\label{p:cases}
(i) A closed emulation is a disembodied person. (ii) An upload whose sensory and motor couplings are severed is a disembodied person, if thinking is functional. (iii) No living human being is disembodied, however complete her paralysis and sensory deprivation. (iv) A brain in a vat connected to stimulating machinery is not disembodied. (v) An immaterial soul with the capacities of a person and no causal link to any body would be a disembodied person, if such a thing is possible.
\end{proposition}

\begin{proof}
(i) is Theorem~3.1 of PID-10-a. (ii) By MND-08 thinking is functional; an upload realizes the functional organization of a reflective thinker (PID-14-a, Assumption~2.3); severed from its couplings, it meets Definition~2.3 of PID-10-a and falls under (i). (iii) A living human being is a thinking organism, or the organism in its personal phase (PID-05), and by Assumption~2.7 a thinking organism perceives through and acts directly on itself, so it is its own body. Paralysis and sensory loss remove channels of perception and action, not the proprioceptive, interoceptive and metabolic states through which the organism registers and regulates itself. (iv) The machinery through which the vat stimulates the brain is something through which the brain-borne person perceives, and so is its body in the sense of Assumption~2.6. (v) follows from the definitions, conditionally on the possibility of such a soul, which nothing here establishes.
\end{proof}

Case (v) is the one \citet[VI]{Descartes1641} took to be actual, and the one that \citet[ch.~3]{Strawson1959} allowed as conceivable only for a \enquote{former person}, whose idea of itself depends on its having once been embodied. The criterion of PID-10-a requires neither dualism nor such a history: case (i) shows that disembodiment is possible for a thinker that never had a body, without a soul.

\begin{corollary}\label{c:human}
No human person becomes a disembodied person; she may have a disembodied successor.
\end{corollary}

\begin{proof}
While a human person exists, she is an organism in its personal phase, and by (iii) she is embodied. A disembodied person arising from her, such as a decoupled upload, is not an organism (PID-10-a, Proposition~3.2), and by PID-14-a, Theorem~3.1, no upload is the original person. By Theorem~3.2 of PID-14-a it is her successor.
\end{proof}

\section{Conclusion}

The possible disembodied persons are computational systems with the organization of a reflective thinker and no coupling to the world, and, conditionally, souls. Human beings in extremity are not among them, and none of us could become one. \enquote{Person} is a phase sortal that can apply to such a system; but the system is a further person, not one of us.

\begin{thebibliography}{99}\small
\bibitem[Descartes(1641)]{Descartes1641} Descartes, R. (1641). \emph{Meditations on First Philosophy}. Cited by meditation.
\bibitem[Strawson(1959)]{Strawson1959} Strawson, P.~F. (1959). \emph{Individuals: An Essay in Descriptive Metaphysics}. Methuen.
\end{thebibliography}

\end{document}
